Supplying designated gold evidence is widely assumed to make language-model fact verification more reliable.
We test that assumption in a paired experiment on 10{,}000 balanced FEVER Supported/Refuted claims, evaluating GPT-5.4 and GPT-5.4-mini without evidence and with FEVER-designated gold evidence (40{,}000 predictions).
Evidence raises aggregate accuracy from 88.69\% to 96.04\% ($+7.35$~pp) for GPT-5.4 and from 84.67\% to 95.54\% ($+10.87$~pp) for GPT-5.4-mini.
Those gains conceal a structured minority of evidence-induced regressions: 226 unique claims that were correct without evidence become wrong after the same gold evidence is supplied (40 shared by both models; 186 specific to one model).
Blinded five-judge progressive-disclosure adjudication on the diagnostic cohort decomposes those regressions into a mixture of evidence-utilization failure, residual warrant ambiguity, and representation-sensitive behavior.
A Cursor/Grok panel assigns 50.0\% / 11.9\% / 38.1\% of the 226 regressions to those categories.
An independent Claude five-judge A$\rightarrow$B$\rightarrow$C replication on the same 226 items (3{,}390 judgments) recovers the same qualitative ordering but not the same prevalences (65.9\% / 5.3\% / 28.8\%; exact taxonomy agreement 74.3\%, $\kappa=0.537$).
External evidence can therefore improve overall reliability while creating localized failures, and automated failure-mode prevalence itself depends partly on the adjudicating model family.
